\documentclass[aspectratio=169]{beamer}
\usetheme{Madrid}
\usepackage[T1]{fontenc}
\usepackage[utf8]{inputenc}
\usepackage{lmodern,booktabs,tabularx,array,tikz}
\usetikzlibrary{arrows.meta,positioning,calc}
\newcolumntype{L}{>{\raggedright\arraybackslash}X}
\newcolumntype{P}[1]{>{\raggedright\arraybackslash}p{#1}}
\setlength{\tabcolsep}{4pt}
\renewcommand{\arraystretch}{1.15}
\title{DepthScope}
\subtitle{Data acquisition and pre-processing}
\author[DepthScope team]{Igor Kołodziej \and Adam Kaniasty \and Mateusz Iwaniuk \and Hubert Kowalski}
\date{22 November 2026}
\begin{document}

\begin{frame}
\titlepage
\end{frame}

\begin{frame}{Why DepthScope?}
\begin{columns}[T,onlytextwidth]
\column{0.58\textwidth}
\small
\textbf{Problem.}\par\vspace{0.35em}
A quoted price does not show how much can be traded near it. For larger BTC or ETH orders, visible depth, spread, fees and venue choice change execution cost.
\medskip

\textbf{Who it is for.}\par\vspace{0.35em}
An analyst comparing liquidity across venues before a trade.
\medskip

\textbf{What we build.}\par\vspace{0.35em}
\begin{itemize}\setlength{\itemsep}{0.3em}
\item Live pipeline for \textbf{Coinbase Exchange} and \textbf{Kraken}: BTC/USD and ETH/USD.
\item Compare visible liquidity and \textbf{forecast} it \textbf{60\,s} ahead (every 5\,s).
\item Research views on execution cost and recorded policies.
\end{itemize}
\column{0.40\textwidth}
\centering
\vspace{0.6em}
\begin{tikzpicture}[
  font=\footnotesize,
  box/.style={draw,structure.fg,rounded corners=2pt,align=center,inner sep=5pt,text width=3.35cm},
  arr/.style={-{Stealth[length=2mm]},structure.fg,line width=0.9pt}
]
  \node[box,fill=structure.fg!12] (ex) {Coinbase and Kraken\\public WebSocket feeds};
  \node[box,below=0.85cm of ex] (ds) {DepthScope\\dashboard and forecasts};
  \node[box,below=0.85cm of ds,fill=structure.fg!12] (an) {Execution analyst\\selects asset and trade side};
  \draw[arr] (ex) -- (ds);
  \draw[arr] (ds) -- (an);
  \draw[arr] (an.west) -- ++(-0.65,0) |- (ds.west);
\end{tikzpicture}
\end{columns}
\end{frame}

\begin{frame}{Data sets and available fields}
\footnotesize
\begin{tabularx}{\linewidth}{@{}P{19mm}P{12mm}LL@{}}
\toprule
\textbf{Data set} & \textbf{Format} & \textbf{Available fields} & \textbf{Planned use}\\
\midrule
Book events & JSON & Price, quantity, side, symbol, exchange time; Kraken checksum. & Depth, spread and imbalance inputs for liquidity regression.\\\addlinespace
Trades & JSON & Trade ID, price, quantity, timestamp and source-side information. & Trade-flow and volume inputs.\\\addlinespace
Metadata & JSON & Trading status, price/quantity increments, heartbeat and connection state. & Data-quality filtering and valid market selection.\\\addlinespace
Fee schedule & CSV & Venue, effective/observed date, maker/taker bps, tier and source URL. & Later execution-cost scenarios.\\
\bottomrule
\end{tabularx}\par
\medskip
Original messages are retained in ClickHouse and Parquet. Normalized acquisition records contain validated book levels and instrument metadata.
\end{frame}

\begin{frame}{Sources and WebSocket channels}
\footnotesize
\begin{tabularx}{\linewidth}{@{}P{17mm}P{26mm}LP{23mm}@{}}
\toprule
\textbf{Venue} & \textbf{Channel} & \textbf{Message content} & \textbf{Arrival}\\
\midrule
Coinbase & \texttt{level2\_batch} & Snapshot and absolute book-level changes. & 50 ms batches when present\\
Coinbase & \texttt{matches} & Executed price, size, trade ID and time. & Event driven\\
Coinbase & \texttt{heartbeat} & Last trade ID and feed continuity information. & Every second\\
Coinbase & \texttt{status} & Market flags and instrument increments. & Periodic\\
\midrule
Kraken v2 & \texttt{book} & Source levels, symbol, time and CRC32; depth 1,000 per side. & Event driven\\
Kraken v2 & \texttt{trade} & Trade ID, side, price, quantity and time. & Event driven\\
Kraken v2 & \texttt{instrument} & Trading status and instrument increments. & Snapshot + updates\\
Kraken v2 & \texttt{heartbeat}, \texttt{status} & Feed and venue health. & Control messages\\
\bottomrule
\end{tabularx}\par\vspace{2mm}
Both feeds cover BTC/USD and ETH/USD. One message may contain several changes or trades.
\end{frame}

\begin{frame}{Data acquisition and storage}
\centering
\includegraphics[width=\linewidth]{source/architecture.pdf}\par\medskip
{\small Fee CSV $\rightarrow$ Python fee loader}
\end{frame}

\begin{frame}{Measured inflow and data quality}
\small
\begin{tabularx}{\linewidth}{@{}Lrrr@{}}
\toprule
\textbf{Venue} & \textbf{Messages} & \textbf{Mean msg/s} & \textbf{Payload kB/s}\\
\midrule
Coinbase Exchange & 4,840 & 40.33 & 30.1\\
Kraken Spot v2 & 31,863 & 265.52 & 51.4\\
\bottomrule
\end{tabularx}\par
\smallskip
{\footnotesize Sample of 120 seconds on 22 November 2026, 18:58:06--19:00:05 UTC. Rates combine BTC, ETH and all subscribed channels.}
\medskip
\begin{itemize}
\item Four books each produced 120 valid records. Nine fixture tests passed.
\item Checks cover absolute quantities, book ordering, Kraken CRC32, Coinbase heartbeats and trading status.
\item Internal trade-ID gaps and a missing initial snapshot need stronger detection. The sample gives a short-window capacity observation.
\end{itemize}
\end{frame}

\begin{frame}{Kappa replay}
\begin{block}{Reprocessing path}
\centering\small
Parquet archive $\rightarrow$ Kafka replay topic $\rightarrow$ C book validation $\rightarrow$ ClickHouse
\end{block}

The replay job reads events in their recorded order and uses the same validator as live acquisition. ClickHouse stores the rebuilt books under a separate run ID.

\begin{block}{Demonstration on a 120-second window}
\begin{itemize}
\item Compared the live and replayed records by venue, asset and second.
\item \textbf{480 of 480} book records matched exactly.
\item Replay processed 43,930 archived events, including warm-up, in 7.4 s.
\end{itemize}
\end{block}
\end{frame}

\begin{frame}{Planned rolling features and forecast timing}
\centering
\input{source/timing.tex}
\vspace{2mm}
\footnotesize
\begin{tabularx}{\linewidth}{@{}P{32mm}P{16mm}L@{}}
\toprule
\textbf{Step} & \textbf{Cadence} & \textbf{Meaning}\\
\midrule
Acquisition (current) & Every 1 s & Save the current validated book and metadata.\\
Features (planned) & Every 1 s & Refresh a window of the last 60 valid samples. For example, sum trade volume or estimate price volatility in that window.\\
Forecasts (planned) & Every 5 s & Use the current features to predict liquidity 60 seconds ahead. Score the saved forecast when its target arrives, then learn.\\
\bottomrule
\end{tabularx}\par
\smallskip
Rolling features summarize this moving window. At $t+1$ s, drop the oldest sample and add the newest. The window keeps 60 samples. The first full window needs 60 consecutive valid acquisition records.
\end{frame}

\begin{frame}{Follow-up work}
\small
\begin{tabularx}{\linewidth}{@{}P{29mm}L@{}}
\toprule
\textbf{Area} & \textbf{Further work}\\
\midrule
Data transformations & Extend source validation and build rolling depth, spread, imbalance, trade-flow and volatility features from the market streams.\\\addlinespace[3mm]
Analytical module & Develop 60-second liquidity regression, compare candidate models with persistence and EWMA, and evaluate them through chronological replay.\\\addlinespace[3mm]
Visualizations & Build Streamlit/Plotly views of venue liquidity, forecasts, observed outcomes and data quality, with asset and trade-side selection.\\
\bottomrule
\end{tabularx}
\end{frame}
\end{document}
